\documentclass{article}
\usepackage{graphicx}
\usepackage{amsmath}
\usepackage{amsfonts}
\usepackage{bookmark}
\usepackage{booktabs}
\usepackage{exercises}
\graphicspath{ {} }
\usepackage{indentfirst}
\usepackage{amsthm}
\usepackage{yhmath}
\usepackage{geometry}
\usepackage{placeins}
\usepackage{csquotes}
\usepackage{pgfplots}
\newcommand*{\oldneg}{\mathord{\sim}}
\newcommand\showdiv[1]{\overline{\smash{)}#1}}
\newcommand\Mydiv[1]{%
\smash{\raise.3ex\hbox{$\big)$}}\mkern-8mu
        \overline{\enspace\strut#1}}
\setlength{\tabcolsep}{5pt}
\theoremstyle{definition}
\newtheorem{definition}{Definition}

\setlength{\parindent}{.5in}
\setlength{\parskip}{12pt}
\renewcommand{\arraystretch}{2}
\geometry{margin=1in}
\newcommand{\?}{\stackrel{?}{=}}
\pgfplotsset{width=7cm,compat=1.18}\usepgfplotslibrary{statistics}

\title{Quiz Answers (pg. 419)}
\begin{document}
\maketitle

\section*{1}
We are given a mean of 47 and a standard deviation of 6. Thus the probability that a radnomly selected $x$-value is between 35 and 65 is 
$13.5+34+34+13.5+2.35=97.35\%$.

\section*{2}
The probability that a randomly selected $x$-value is at least 41 is $34+34+13.5+2.35+0.15=84\%$.

\section*{3}
The probability that a randomly selected $x$-value is at most 29 is $0.15\%$.

\section*{4}
\[ \pm 0.03=\pm\frac{1}{\sqrt{n}} \] \[ 0.0009=\frac{1}{n} \] \[ n\approx 1111 \]

\section*{5}
\[ \pm 0.07=\pm\frac{1}{\sqrt{n}} \] \[ 0.0049=\frac{1}{n} \] \[ n\approx 204 \]

\section*{6}
\[ \pm 0.045=\pm\frac{1}{\sqrt{n}} \] \[ 0.045^2=\frac{1}{n} \] \[ n\approx 494 \]

\section*{7}
\[ \pm 0.008=\pm\frac{1}{\sqrt{n}} \] \[ 0.000064=\frac{1}{n} \] \[ n=15625 \]

\section*{8}
First, it is not always possible to do an experiment in the first place. It could be that experimentation would be unethical, or simply impracticle. But when an experiment is 
is possible, it is almost always better to do a comparitive experiment rather than not. In comparing a control group to a treatment group, we can definitively see exactly 
how much treatment is being applied, and can control the conditions of both groups, which can help validate our results. If we were instead doing an observational study 
(or a non-comparative experiment), it 
might be harder to conclude a causal relationship, and we may only be able to find a correlation. Lastly, it is crucial that our experiments are based upon 
random samples as much as possible. If 
they are not, we risk getting a result which is biased or unrepresentative (of the entire population), which would invalidate any of the findings of the experiment. 

\end{document}
